\documentclass[12pt,a4paper]{article}
\usepackage{cmap} % Makes the PDF copiable. See http://tex.stackexchange.com/a/64198/25761
\usepackage[T1]{fontenc}
\usepackage[brazil]{babel}
\usepackage[utf8]{inputenc}
\usepackage{amsmath}
\usepackage{amsfonts}
\usepackage{amssymb}
\usepackage{amsthm}
\usepackage[usenames,svgnames,dvipsnames]{xcolor}
\usepackage{hyperref}
\usepackage{multicol}
\usepackage{graphicx}
\usepackage[top=2cm, bottom=2cm, left=2cm, right=2cm]{geometry}

\hypersetup{
    colorlinks = true,
    allcolors = {blue}
}

% TODO: Consider using exsheets
% http://linorg.usp.br/CTAN/macros/latex/contrib/exsheets/exsheets_en.pdf
%
% http://ctan.org/tex-archive/macros/latex/contrib/exercise/
% Options: answerdelayed,lastexercise,noanswer
\usepackage[answerdelayed,lastexercise]{exercise}

\addto\captionsbrazil{%
\def\listexercisename{Lista de exerc\'icios}%
\def\ExerciseName{Exerc\'icio}%
\def\AnswerName{Solu\c{c}\~ao do exerc\'icio}%
\def\ExerciseListName{Ex.}%
\def\AnswerListName{Solu\c{c}\~ao}%
\def\ExePartName{Parte}%
\def\ArticleOf{de\ }%
}

\renewcommand{\ExerciseHeaderTitle}{(\ExerciseTitle)\ }
\renewcommand{\ExerciseListHeader}{%\ExerciseHeaderDifficulty%
\textbf{%\ExerciseListName\
\ExerciseHeaderNB.\ %
%\ --- \
\ExerciseHeaderTitle}%
%\ExerciseHeaderOrigin
\ignorespaces}
\renewcommand{\AnswerListHeader}{\textbf{\ExerciseHeaderNB.\ (\AnswerListName)\ }}

\newcommand*\diff{\mathop{}\!\mathrm{d}}
\newcommand*\sen{\operatorname{sen}}
\newcommand*\tg{\operatorname{tg}}
\newcommand*\R{\mathbb{R}}

\renewcommand{\theenumi}{\alph{enumi}}
\renewcommand\labelenumi{(\theenumi) }

\newcommand*\tipo{Prova IV}
\newcommand*\turma{NEXM241-D}
\newcommand*\disciplina{CDI1001}
\newcommand*\eu{Helder G. G. de Lima}
\newcommand*\data{03/07/2024}

\author{\eu}
\title{\tipo - \disciplina}
\date{\data}

\begin{document}
\thispagestyle{empty}
\newgeometry{margin=2cm,bottom=0.5cm}
\begin{center}
\includegraphics[width=9.0cm]{marca} \\
\textbf{\tipo\ (\disciplina / \turma)} \\
Prof. \eu\footnote{
Este é um material de acesso livre distribuído sob os termos da licença \href{https://creativecommons.org/licenses/by-sa/4.0/deed.pt_BR}{Creative Commons BY-SA 4.0}}
\end{center}

\noindent Nome do(a) aluno(a): \underline{\hspace{9,7cm}} Data: \underline{\data}

%\section*{Instruções}
\begin{center}\fbox{
\begin{minipage}{14cm}

{\footnotesize
\begin{itemize}
\renewcommand{\theenumi}{\Roman{enumi}}
\item Identifique-se em todas as folhas.
\item Mantenha o celular e os demais equipamentos eletrônicos desligados durante a prova.
\item Justifique cada resposta com cálculos ou argumentos baseados na teoria estudada.
\item Resolva $4$ das $5$ questões (deixe claro que questão não deverá ser corrigida).
\end{itemize}
}

\end{minipage}
}
\end{center}

%\section*{Questões}
\begin{ExerciseList}
\Exercise[title={2,0}] Para cada item abaixo, diga se é verdadeiro (V) ou falso (F).

\emph{Observação 1:} Não é necessário justificar.

\emph{Observação 2:} Cada item marcado errado anula a pontuação de um item marcado corretamente. Você tem a opção de deixar o item em branco e, nesse caso, você nem ganha e nem perde pontos.

\begin{enumerate}
    \item $\boldsymbol{\left(\quad\right)}$ Se $f(x)$ tem concavidade para cima no intervalo $(a, b)$ então ela é crescente em $(a, b)$.
    \item $\boldsymbol{\left(\quad\right)}$ Toda assíntota horizontal de uma função $g(x)$ é uma assíntota oblíqua de $g(x)$.
    \item $\boldsymbol{\left(\quad\right)}$ Se $g(x)$ e $h(x)$ são primitivas de $f(x)$ em $\R$ então existe $c \in \R$ tal que $g(x) = h(x) + c$.
    \item $\boldsymbol{\left(\quad\right)}$ Para quaisquer $f(x)$ e $g(x)$, vale a igualdade $\int f(x) g(x) \diff{x} = \int f(x) \diff{x} \cdot \int g(x) \diff{x}$.
\end{enumerate}

\Answer \begin{enumerate}
    \item \textbf{Falso}: Existem funções decrescentes com concavidade para cima, tais como $f(x) = e^{-x}$, para $x \in \R$ e $g(x) = \frac{1}{x}$ para $x \in (0, +\infty)$)
    \item \textbf{Verdadeiro}: Se $y=k$ é uma assíntota horizontal de $g(x)$, então $\lim_{x\to\infty} g(x) = k$. Logo, $\lim_{x\to \infty} \frac{g(x)}{x} = 0$ e $\lim_{x\to \infty} g(x) - 0 x = k$, ou seja $y=0x + k$ é uma assíntota oblíqua.
    \item \textbf{Verdadeiro}: Se $g(x)$ e $h(x)$ são primitivas de $f(x)$ em $\R$, então $g^\prime(x) = h^\prime(x)$, e portanto $(g(x) - h(x))^\prime = g^\prime(x) - h^\prime(x) = 0$. Logo, existe $c \in \R$ tal que $g(x) - h(x) = c$, isto é, $g(x) = h(x) + c$.
    \item \textbf{Falso}: Por exemplo, se $f(x) = 2$ e $g(x) = 3$ $\int f(x) g(x) \diff{x} = \int 2\cdot 3 \diff{x} = 6x + c$ mas $\int f(x) \diff{x} \cdot \int g(x) \diff{x} = \int 2 \diff{x} \cdot \int 3 \diff{x} = 2x \cdot 3x + c = 6x^2 + c \neq 6x + c$.
\end{enumerate}

\Exercise[title={2,0}] Um fazendeiro tem $120m$ de cerca para construir dois galpões retangulares adjacentes (de mesma área), conforme a figura a seguir, em que as linhas tracejadas indicam onde haverá cerca. Quais devem ser as dimensões totais dos galpões de modo a maximizar sua área total?

\begin{center}
    \includegraphics[width=7cm]{img/prova-4-nex-galpões.pdf}
\end{center}

\Answer Sejam $x > 0$ e $y > 0$ as medidas (na horizontal e na vertical, respectivamente) de um dos galpões retangulares.
Pela disposição dos galpões, e pela quantidade de cerca disponível, deve ocorrer $4x + 3y = 120$. Consequentemente:
\[
    y = \frac{1}{3}(120 - 4x) = 40 - \frac{4}{3}x < 40
\]
e
\[
    x = \frac{1}{4}(120 - 3y) = 30 - \frac{3}{4}y < 30
\]
Assim, a área total é dada por $A(x) = 2x\cdot y = 2x \cdot (40 - \frac{4}{3}x) = 80x - \frac{8}{3}x^2$, em que $0 < x < 30$.

O valor de $x$ que maximiza a área total deve ser um ponto crítico de $A(x)$, e tem-se:
\[
    A^\prime(x) = 80 -\frac{16}{3}x, \text{ para todo } x \in (0, 30).
\]
Neste caso,
\[
    A^\prime(x) = 0
    \Leftrightarrow
    80 -\frac{16}{3}x = 0
    \Leftrightarrow
    80 = \frac{16}{3}x
    \Leftrightarrow
    240 = 16x
    \Leftrightarrow
    x = \frac{240}{16} = 15.
\]

Como $A^{\prime\prime}(x) = -\frac{16}{3}$, para todo $x \in (0, 30)$, e em particular $A^{\prime\prime}(15) < 0$, segue do teste da segunda derivada que $x = 15$ é um ponto de máximo de $A(x)$. O valor de $y$ correspondente é $y = 40 - \frac{4}{3} \cdot 15 = 20$. Portanto, os galpões devem medir $15m$ por $20m$ para que a área total seja máxima.

\Exercise[title={2,0}] Encontre uma primitiva $F$ da função $f(x) = x \sen(x)$ que satisfaça $F(\pi/2) = 0$.
\Answer A integral indefinida de $f$ fornece todas as primitivas de $f$. Usando a integração por partes com $u = x$, $\diff{v} = \sen(x) \diff{x}$, tem-se $\diff{u} = \diff{x}$, $v = -\cos(x)$ e

\begin{align*}
    \int x \sen(x) \diff{x}
    & = x (-\cos(x)) - \int -\cos(x) \diff{x} \\
    & = - x \cos(x) + \int \cos(x) \diff{x}
      = - x \cos(x) + \sen(x) + c.
\end{align*}

Para que $F(x) = - x \cos(x) + \sen(x) + c$ satisfaça $F(\pi/2) = 0$ é preciso que
\[
0
= F(\pi/2)
= F(\pi/2)
= - \pi/2\cdot \cos(\pi/2) + \sen(\pi/2) + c.
= 0 + 1 + c
\]
Logo, $c = -1$ e $F(x) = - x \cos(x) + \sen(x) -1$.

\Exercise[title={2,0}] Calcule as integrais
\begin{multicols}{2}
    \begin{enumerate}
        \item $\displaystyle\int \frac{\sqrt{x}}{2x^2} - \frac{2x^2}{\sqrt{x}} \diff{x}$
        \item $\displaystyle\int \sen(2024 x) \cos(2023 x) \diff{x}$
    \end{enumerate}
\end{multicols}
\Answer
\begin{enumerate}
\item
\begin{align*}
    \int \frac{\sqrt{x}}{2x^2} - \frac{2x^2}{\sqrt{x}} \diff{x}
    & = \int \frac{1}{2}x^{\frac{1}{2} - 2} - 2 x^{2-\frac{1}{2}} \diff{x}
      = \frac{1}{2} \int x^{-\frac{3}{2}} \diff{x} - 2 \int x^{\frac{3}{2}} \diff{x} \\
    & = \frac{1}{2} \frac{x^{-\frac{1}{2}}}{(-\frac{1}{2})} - 2 \frac{x^{\frac{5}{2}}}{\frac{5}{2}} + c
      = - x^{-\frac{1}{2}} - \frac{4}{5} x^{\frac{5}{2}} + c
      = -\frac{1}{\sqrt{x}} - \frac{4}{5} \sqrt{x^5} + c
\end{align*}

\item
\begin{align*}
    \int \sen(2024 x) \cos(2023 x) \diff{x}
    & = \int \frac{1}{2}\left[ \sen(2024 x - 2023 x) + \sen(2024 x + 2023 x)\right] \diff{x} \\
    & = \frac{1}{2} \int \sen(x) \diff{x} + \frac{1}{2} \int \sen(4047 x) \diff{x} \\
    & = \frac{1}{2} (-\cos(x))  + \frac{1}{2} \int \frac{1}{4047} \sen(4047 x) \cdot 4047 \diff{x} \\
    & \quad \quad (u = 4047 x\Rightarrow \diff{u} = 2047 \diff{x}) \\
    & = -\frac{\cos(x)}{2}  + \frac{1}{8094} \int \sen(u) \diff{u} \\
    & = -\frac{\cos(x)}{2} -\frac{1}{8094} \cos(u) + c
      = -\frac{\cos(x)}{2} -\frac{\cos(4047 x)}{8094} + c
\end{align*}
\end{enumerate}

\Exercise[title={2,0}] Calcule as integrais
\begin{multicols}{2}
    \begin{enumerate}
        \item $\displaystyle\int \frac{\sec(\ln(x))\tg(\ln(x))}{x} \diff{x}$
        \item $\displaystyle\int \dfrac{12x^3}{e^{3x^4 - 45}} \diff{x}$
    \end{enumerate}
\end{multicols}
\Answer
\begin{enumerate}
    \item
    \begin{align*}
        \int \frac{\sec(\ln(x))\tg(\ln(x))}{x} \diff{x}
        & = \int \sec(\ln(x))\tg(\ln(x)) \cdot \frac{1}{x} \diff{x} \\
        & \quad \quad (u = \ln(x)\Rightarrow \diff{u} = \frac{1}{x} \diff{x} ) \\
        & = \int \sec(u)\tg(u) \diff{u}
          = \sec(u) + c
          = \sec(\ln(x)) + c
    \end{align*}

    \item
    \begin{align*}
        \int \dfrac{12x^3}{e^{3x^4 - 45}} \diff{x}
        & = \int \dfrac{1}{e^{3x^4 - 45}} \cdot 12x^3 \diff{x} \\
        & \quad \quad (u = 3x^4 - 45 \Rightarrow \diff{u} = 12x^3 \diff{x} ) \\
        & = \int \dfrac{1}{e^{u}} \diff{u}
          = \int e^{-u} \diff{u}
          = -\int e^{-u} \cdot (-1) \diff{u} \\
        & \quad \quad (v = -u \Rightarrow \diff{v} = (-1) \diff{u} ) \\
        & = -\int e^{v} \diff{v}
          = -e^{v} + c
          = -e^{-u} + c
          = -e^{-(3x^4 - 45)} + c
          = -e^{45 - 3x^4} + c.
    \end{align*}
    \end{enumerate}
\end{ExerciseList}

\begin{center}
BOA PROVA E BOAS FÉRIAS!
\end{center}

\newpage
\restoregeometry
\section*{Respostas}
\shipoutAnswer
\end{document}
